\chapter{A local filling condition from subtree polynomials}
\label{ch:local-circle-messages}

The summed condition in Chapter~\ref{ch:quantitative-circle-trees}
compares an error for the complete augmented diagram with the total
change of action caused by filling. That comparison is useful when
the sum of the changes is small. A tree also permits a different
calculation. Evaluate a terminal subtree first, keep its value at
zero exactly, and estimate the polynomial left on its one boundary.
The next trace then needs to absorb only its own change of action.
This chapter proves that this procedure gives a sufficient condition
at each vertex, with no sum over the vertices.

We keep exactly the ordered tangle construction of
Section~\ref{sec:circle-the-knot-construction}. There is one crossing
circle at every vertex of a finite tree, including its terminal
vertices. The cyclic orders, positions of the crossing boxes,
quarter-turns, disk framings and retained half twists are part of
the diagram. Each tree edge is an actual quarter-turn. Write
$d_v\ge1$ for a vertex degree and
$n_v=\epsilon_v-2p_v$ for its signed crossing power. The component
test must give one main knot, including no closed main component
inside an intermediate tangle.

\section{The local volume theorem}

The reduced invariant, the root of unity, and the exact scalar
factors have the same normalization as before:
\begin{equation}\label{eq:local-normalizations}
\begin{aligned}
 q_N&=e^{\pi i/N},& K_N(K)&=F_N(1),\\
 H_N(s)&=\frac{\sin\pi s}{N\sin(\pi s/N)},&
 A_N(s)&=q_N^{a(D)s^2/2},\\
 a(D)&=-w(D)-\sum_vn_v.
\end{aligned}
\end{equation}
The removable value of $H_N$ at zero is one. The exponent $a(D)$
uses the oriented writhe of the actual filled diagram and its
original checked-crossing convention.

\begin{theorem}[A local degree and twist condition]
\label{thm:local-circle-volume}
Fix the ordered tree, its strand frames, its knot parity pattern,
and all its signed integer powers. Suppose
\begin{equation}\label{eq:local-criterion}
 \kappa=\max_v\frac{d_v}{d_v^2+n_v^2}\le2^{-120}.
\end{equation}
Then the actual filled knot complement is hyperbolic. For every
fixed $0\le R<1$,
\begin{equation}\label{eq:local-profile}
 \frac{F_N(q_N^s)}{F_N(1)A_N(s)H_N(s)}
       =1+O_{\mathcal T,\mathbf n,R}(N^{-1}),
 \qquad |s|\le R,
\end{equation}
and
\begin{equation}\label{eq:local-volume}
 \frac{2\pi}{N}\log|F_N(q_N^s)|
       \longrightarrow\Vol(S^3\setminus K)
\end{equation}
uniformly on that disk. These values are nonzero for all sufficiently
large levels. The limit includes both parities and composite levels.
The tree, the powers and the disk are fixed before $N$ tends to
infinity.
\end{theorem}

The numerical test is independent of the number of vertices.
Polynomial prefactors, error constants and starting levels can
depend on the entire fixed tree and its powers. Thus the theorem
does not assert a limit for a sequence of diagrams growing with $N$.
The knot condition also retains the two-circle qualification from
Chapter~\ref{ch:quantitative-circle-trees}: the two-odd pattern is
not a knot, so an admissible two-circle knot has a pure leaf.

The constants in the summed and local tests are different. Neither
test is being used as an optimal characterization. Both remain
sufficient results. A simple consequence of the local test is
\[
 \frac{d}{d^2+n^2}\le\frac1{2|n|}\quad(n\ne0):
 \qquad |n_v|\ge2^{119}\ \text{for every }v
 \quad\Longrightarrow\quad\eqref{eq:local-criterion}.
\]
This bound no longer grows with the number of edges. A zero power
is permitted if its degree is at least $2^{120}$, provided the
other vertices satisfy their own inequalities and the diagram is
a knot. These large constants make the scope explicit; they are
not estimates of the smallest admissible twists.

\section{What a subtree leaves on its boundary}

Choose a terminal root. A nonroot vertex has one edge toward its
parent and a finite ordered list of child subtrees. After evaluating
all the crossings and internal edges in those child subtrees,
each child leaves a polynomial in the one variable at its boundary.
Writing that variable as $-2\cos t$ is convenient for the exact
root-of-unity trace. We denote the polynomial value by $M_{w,N}(t)$.
It is even and periodic as a function of the angle $t$.

The fixed Newton polynomials $e_k$ and their full pairing are
defined in Section~\ref{sec:q-coordinates}. The full edge kernel is
\begin{equation}\label{eq:local-kernel}
 K_N(t,u;s)=\sum_{k=0}^{N-1}
 e_k(-2\cos t)e_k(-2\cos u)\frac{H_N(s)}{W_k(s)}.
\end{equation}
Every Newton index is included. In particular the boundary
polynomial is not obtained by discarding high indices before the
subtree has been evaluated.

Let $\mathcal D_{n_v,\sigma_v,N}^{(t)}$ be the actual trace after
multiplication by the full crossing power at $v$. It acts in
$-2\cos t$, takes the remainder in the prescribed Casimir quotient,
and retains values and first derivatives at repeated roots.
The orientation type determines $\sigma_v\in\{0,s,-s\}$. The edge
orientation scalar $\gamma_{e_v}$ belongs to that edge once.
The exact outgoing polynomial is
\begin{equation}\label{eq:local-recursion}
 M_{v,N}(u)=\gamma_{e_v}(s)\,
 \mathcal D_{n_v,\sigma_v,N}^{(t)}
 \left[K_N(t,u;s)\prod_{w\text{ child of }v}M_{w,N}(t)\right].
\end{equation}
Only $K_N$ depends on $u$. Hence $M_{v,N}(u)$ has trigonometric
degree at most $N-1$, independently of the size or depth of the
subtree. The trace may receive a polynomial of much larger degree
in $-2\cos t$; its quotient operation handles that degree without
changing the outgoing degree bound.

\begin{figure}[ht]
\centering
% The diagram is the actual finite tensor contraction on a rooted tree.
\begin{tikzpicture}[x=1cm,y=1cm,line cap=round,line join=round,
  font=\small,
  vertex/.style={circle,draw=blue!65!black,fill=blue!8,
    minimum size=8mm,inner sep=1pt},
  edge/.style={line width=1pt,draw=blue!65!black},
  sub/.style={draw=gray!65,fill=gray!7,line width=.7pt}]
\node[vertex] (parent) at (0,3.15) {$u$};
\node[above=5pt] at (parent.north) {Parent angle};
\node[vertex] (v) at (0,.8) {$t$};
\draw[edge] (parent)--(v);
\node[fill=white,inner sep=3pt,align=center] at (-.9,2.0)
 {$K_N(t,u;s)$};
\node[right=8pt,align=left] at (v.east)
 {Vertex $v$\\crossing power $n_v$};
\coordinate (left) at (-1.8,-.65);
\coordinate (middle) at (0,-.65);
\coordinate (right) at (1.8,-.65);
\foreach \p in {left,middle,right}{
  \draw[edge] (v)--(\p);
  \path[sub] (\p)--($(\p)+(-.55,-1.25)$)--
    ($(\p)+(.55,-1.25)$)--cycle;
}
\node at (-1.8,-1.5) {$w_1$};
\node at (0,-1.5) {$\cdots$};
\node at (1.8,-1.5) {$w_m$};
\node[fill=white,inner sep=2pt,rotate=39] at (-1.3,-.15)
 {$M_{w_1,N}(t)$};
\node[fill=white,inner sep=2pt,rotate=-39] at (1.3,-.15)
 {$M_{w_m,N}(t)$};
\node[align=center,text width=5.2cm] at (0,-3.05)
 {Each child subtree has already been evaluated\\
 as a polynomial in $-2\cos t$.};
\node[anchor=north west,align=left,text width=7.4cm] (kerneltext)
 at (3.3,3.65) {
 The edge keeps every index $0\le k<N$:
 \[
 K_N(t,u;s)=\sum_{k=0}^{N-1}
 e_k(-2\cos t)e_k(-2\cos u)\frac{H_N}{W_k}.
 \]
 The trace at $v$ acts only in $t$:
 \[
 M_{v,N}(u)=\gamma_{e_v}\,
 \mathcal D_v^{(t)}
 \left[K_N(t,u;s)\prod_{j=1}^{m}M_{w_j,N}(t)\right].
 \]
 Trigonometric degree in $u$: at most $N-1$.
 };
\node[anchor=north west,align=left,text width=7.4cm]
 at ([yshift=-4mm]kerneltext.south west) {
 Factor out the exact value at zero:
 \[
 \widehat M_{v,N}(u)=\frac{M_{v,N}(u)}{M_{v,N}(0)}.
 \]
 Its local phase and its real decay control the next
 trace. The denominator includes all crossing and
 orientation factors from this subtree.
 };
\end{tikzpicture}

\caption{Evaluating a subtree of the original finite contraction.
The blue edge retains the full Newton kernel. Each lower triangle
denotes a child subtree with its existing polynomial boundary value.
The trace at $v$ acts only on $t$, leaving an even trigonometric
polynomial in the parent angle $u$. Its value at zero is factored
out exactly before estimating the next trace.}
\label{fig:tree-boundary-message}
\end{figure}

The last root trace acts on the product of its child polynomials.
The original closing and framing factors, reference factors and
one reduced cut factor are then restored. This is a change in the
order of a finite tensor contraction. It does not commute tangles
or change any physical quarter-turn.

\section{Why the normalization is local}

The induction proves that $M_{v,N}(0)$ is nonzero and that, near
the origin,
\begin{equation}\label{eq:local-message-expansion}
 \frac{M_{v,N}(u)}{M_{v,N}(0)}
 =\frac{\alpha_v(u)}{\alpha_v(0)}
   e^{Nh_v(u)}(1+O(N^{-1})).
\end{equation}
Here $\alpha_v$ is holomorphic and nonzero,
$h_v(0)=h'_v(0)=0$, and $h_v$ is even. The important uniform
property is its curvature:
\begin{equation}\label{eq:local-phase-class}
 \left|h_v''(u)+\frac1{2\pi}\right|
       \le\frac1{40\pi},\qquad |u|\le2^{-12}.
\end{equation}
It depends on the single edge joining the subtree to its parent.
It does not grow with the number of vertices already evaluated.

To see why this is possible, let $\psi(t,u)$ be the exact edge
action of \eqref{eq:qe-9}--\eqref{eq:qe-12}. The conditional action
at a vertex is
\[
 A_v(t,u)=\psi(t,u)+\sum_w h_w(t).
\]
Its $t$ curvature is close to $-d_v/(2\pi)$. After including the
crossing chirp and the first Fourier charge, the stationary
equation has the form
\[
 \left(\frac{d_v-in_v}{2\pi}\right)t=-i+E_{v,t}(t,u),
 \qquad |E_{v,tt}|\le\frac{d_v}{40\pi}.
\]
The stationary value, with its value at $u=0$ subtracted, is the
new phase $h_v$. Differentiating a stationary value gives a Schur
complement:
\[
 h_v''=\psi_{uu}
 -\frac{\psi_{ut}^2}{A_{v,tt}+in_v/(2\pi)}.
\]
The denominator grows with $d_v$ and $|n_v|$, whereas the two
derivatives in the numerator come from one edge. The explicit
edge bounds therefore preserve \eqref{eq:local-phase-class}.

The exact denominator in \eqref{eq:local-message-expansion} is
essential. On a smaller zero-free disk,
\[
 \frac1N\log\frac{M_{v,N}(u)}{M_{v,N}(0)}
 =h_v(u)+\frac1N\log\frac{\alpha_v(u)}{\alpha_v(0)}
       +O(N^{-2}).
\]
The next trace retains this entire expression. Its error changes
an analytic amplitude; it does not add the difference between a
subtree's complete and filled actions to each of its ancestors.

\section{From a local saddle to the full finite trace}

A local expansion alone cannot exclude a larger contribution far
from the origin. The induction also proves a real angular envelope
for the exact polynomial,
\[
 \left|\frac{M_{v,N}(u)}{M_{v,N}(0)}\right|
 \le CN^A e^{-N2^{-18}(1-\cos u)},\qquad u\in[-\pi,\pi].
\]
The degree bound following \eqref{eq:local-recursion} converts this
to a linear bound on complex strips by the subharmonic maximum
principle. Near the origin the phase curvature gives a quadratic
bound instead. Combining the two proves one fixed complex-strip
estimate for every subtree.

The exact crossing trace is then treated on its translated root
grid. Its Fourier coefficients retain the Chebyshev factors
$U_{\ell-1}(\cos2\pi\sigma_v)$, the parity signs and all endpoint
currents. The first coefficient is exactly one. The conditional
critical values are strictly ordered with the charge, and all
large charges and contour connectors have a fixed conditional
loss. Condition~\eqref{eq:local-criterion} makes the current
vertex's first-charge action drop smaller than that loss.
It does not compare this loss with a sum over descendants.

Finally the real envelope is proved afresh from the actual trace:
the full polynomial norm of a fixed crossing power, multiplication,
quotient reduction and the closing functional cost only a polynomial
in $N$. The edge kernel supplies angular decay at the outgoing
variable. Dividing by the newly proved nonzero value at zero spends
only the same local action drop. This closes the induction.
Appendix~\ref{app:local-messages} gives all constants, finite
contours, collision limits and normalizations.

\section{A separate vertexwise metric criterion}

The geometric completion requires less than the local quantum
condition. The following result is useful independently.

\begin{theorem}[Positive completion from vertexwise data]
\label{thm:vertexwise-geometry}
For the same actual ordered construction, suppose
\begin{equation}\label{eq:vertexwise-criterion}
 L_v=\sqrt{d_v^2+n_v^2}\ge630\qquad\text{for every }v.
\end{equation}
Then the charge equations
\begin{equation}\label{eq:vertexwise-charge}
 \Psi_v(\boldsymbol\tau)+\frac{in_v}{2\pi}\tau_v=ic,
 \qquad \Psi=\sum_{vw}\psi(\tau_v,\tau_w),\qquad 0\le c\le1
\end{equation}
have a unique solution in $|\tau_v|\le10^{-2}$. It is holomorphic
in $c$ near the closed unit disk and satisfies
\begin{equation}\label{eq:vertexwise-size-sign}
 |\tau_v(c)|\le
 \frac{2\pi c}{L_v-d_v/39000},\qquad
 \Im\tau_v(c)<0\quad(c>0).
\end{equation}
The actual cells complete to cone solid tori of angle $2\pi c$,
and to smooth solid tori at $c=1$. The main cusp remains complete.
When the component test gives a knot, the endpoint is its complete
finite-volume hyperbolic complement.
\end{theorem}

The sign in \eqref{eq:vertexwise-size-sign} is proved, rather than
inferred from the smallness of the parameters. Put
\[
 F(t,u)=2\operatorname{arsinh}(\tan(t/2)\cos(u/2)),\qquad E=F-t.
\]
On the radius $10^{-2}$ bidisk, the explicit derivatives give
$|E_t|,|E_u|<1/39000$. Since $E(0,u)=0$, write $E=tH$
with $H$ holomorphic and $|H|<1/39000$. The charge equation is
\[
 \left(d_v-in_v+\sum_{w\sim v}H(\tau_v,\tau_w)\right)\tau_v
       =-2\pi ic.
\]
Its denominator has strictly positive real part. This gives the
sign and a contraction whose row bound is $2/39000$, independent
of the tree size. The proof, including a finite variant with
different radii at different vertices, is in
Section~\ref{sec:vg-derivatives}.

The original primitive periods are
\begin{equation}\label{eq:vertexwise-periods}
 U_v=-2\pi\Psi_v-i\epsilon_v\tau_v,\qquad
 V_v=2i\tau_v,\qquad U_v+p_vV_v=-2\pi ic.
\end{equation}
Appendix~\ref{app:quantitative-geometry} proves that these are
actual logarithmic periods along a positive path. At the initial
cusp their integral ambiguities vanish; continuity preserves that
choice. At $c=1$ the primitive filling slope therefore makes one
full turn. The conclusion concerns the original marked exterior,
even if a total meridian logarithm has moved far from zero.

\section{The original volume and the two sufficient tests}

Back-substitution of the subtree stationary points gives the
simultaneous equations \eqref{eq:vertexwise-charge} at $c=1$.
The constants in the subtree phases cancel exactly, leaving
\[
 \mathcal C(c)=\operatorname*{stat}_{\boldsymbol\tau}
 \left[\Psi+\frac{i}{4\pi}\sum_vn_v\tau_v^2
                    -ic\sum_v\tau_v\right].
\]
The primitive-period calibration in
Section~\ref{sec:fg-6} gives
$\Vol_c=2\pi\Re\mathcal C(c)$. This identifies the action selected
by the finite trace with the volume of the actual filled complement.
It is the geometric step in Theorem~\ref{thm:local-circle-volume}.

Each edge sum contributes $N^{-3/2}$, and each circle trace contributes
$N^2$ times a nondegenerate one-dimensional Gaussian. A tree has one
more vertex than edge, so the full leading power is $N^{3/2}$.
Its amplitude is nonzero. The exact first Fourier coefficient,
the one reduced cut factor and the original scalar ledger then
give \eqref{eq:local-profile} and \eqref{eq:local-volume}.

The geometric criterion~\eqref{eq:vertexwise-criterion} by itself
does not prove a quantum volume limit. It can hold when neither
the summed quantum condition nor the local quantum condition does.
For example, on arbitrarily long paths with powers chosen from
$630$ and $631$, it holds at every vertex while the sum
$\sum_vd_v/(d_v^2+n_v^2)$ grows with the number of vertices.
The appropriate knot parity still has to be selected.
Theorem~\ref{thm:local-circle-volume} supplies a separate analytic
criterion. The earlier summed criterion in
Theorem~\ref{thm:circle-quantitative} remains useful, often with
much smaller powers on a particular finite diagram.
